\section{What rank-one convexity gives}\label{sec:roc}

Throughout this section $\Phi$ is a profile of degree $p>0$. The section collects what the
constructions use from \ref{P2}. The proofs are in \cref{app:roc}.

\subsection{Two special lines}

\begin{lemma}\label{lem:basic}
\begin{enumerate}[label=\textup{(\alph*)}]
\item The function $\tilde h(s)=\Phi\bigl(\abs{1+s}/2,\abs{1-s}/2\bigr)$ is convex on $\R$;
  it equals $h$ on $[-1,1]$.
\item If $0\le x_0\le x_1$, $0\le y_0\le y_1$ and $x_1^2-x_0^2=y_1^2-y_0^2$, then
  $\Phi(x_0,y_0)\le\Phi(x_1,y_1)$.
\end{enumerate}
\end{lemma}

Part (a) is convexity along the line $s\mapsto\diag(1,s)$. Part (b) says that $\Phi$ grows along
the hyperbolas $x^2-y^2=\mathrm{const}$, on which the determinant is constant; it is the
Baker--Ericksen inequality in integrated form, and it comes from the line
$\tau\mapsto(x_0+i\tau,y_0+i\tau)$, along which the profile is convex and even.

\subsection{Burkholder-type minorants}

\begin{definition}\label{def:minorant}
For $c_0,m\in\R$ the \emph{Burkholder-type function} with the coefficients $(c_0,m)$ is
\[
  G_{c_0,m}(x,y)=(x+y)^{p-1}\bigl(c_0(x+y)+m(x-y)\bigr),\qquad x,y\ge0 .
\]
\end{definition}

It is the profile of $A\mapsto\abs A^{p-2}(c_0\abs A^2+m\det A)$, and
$G_{c_0,m}(x,y)=\sigma^p\ell(\delta/\sigma)$ with the affine function $\ell(s)=c_0+ms$: the
Burkholder-type functions are the homogeneous functions that are affine on the segment
$x+y=1$. The Burkholder function is $B_p=G_{\abs\kappa,1}$ with $\kappa=(p-2)/p$. Hence
$G_{c_0,m}\le\Phi$ on the quadrant if and only if $\ell\le h$ on $[-1,1]$. Since $\tilde h$ is
convex on $\R$, $h$ has a supporting line at every $t\in[-1,1]$:

\begin{lemma}\label{lem:minorant}
For every $t\in[-1,1]$ there is $m\in\R$ such that, with $H=h(t)$, the Burkholder-type function
$G_{H-mt,m}$ lies below $\Phi$ on the quadrant and equals $\Phi$ at $(x_t,y_t)$.
\end{lemma}

\subsection{The first-order inequality}

A convex function lies above its tangent lines. For a profile, which need not be differentiable,
the minorant supplies the tangent.

\begin{proposition}[The first-order inequality]\label{prop:tangent}
Let $\abs t<1$, $H,m\in\R$, and suppose that $G_{H-mt,m}\le\Phi$ on the quadrant and
$\Phi(x_t,y_t)=H$. Put
\[
  \Phi_x=pH+(1-t)m,\qquad\Phi_y=pH-(1+t)m .
\]
Then for all $\alpha,\beta\in\C$ with $\abs\alpha=\abs\beta$ and all $\tau\in\R$,
\begin{equation}\label{eq:tangent}
  \Phi\bigl(\abs{x_t+\tau\alpha},\abs{y_t+\tau\beta}\bigr)\ge
  H+\tau\bigl(\Phi_x\re\alpha+\Phi_y\re\beta\bigr).
\end{equation}
\end{proposition}

The numbers $\Phi_x$ and $\Phi_y$ are the partial derivatives at $(x_t,y_t)$ of the minorant
$G_{H-mt,m}$; they are those of $\Phi$ where $\Phi$ is differentiable. The proof compares the
convex function $\tau\mapsto\Phi(\abs{x_t+\tau\alpha},\abs{y_t+\tau\beta})$ with the same
function of the minorant, which lies below it, agrees with it at $\tau=0$ and is differentiable
there; a convex function lies above the tangent of such a comparison function
(\cref{lem:convex}).

The inequality \eqref{eq:tangent} holds for every real $\tau$, and the constructions use it far
from $\tau=0$, on two lines. The first joins the base matrix to a rotated copy of it, where
$\Phi$ has the same value $H$:

\begin{corollary}[Baker--Ericksen]\label{cor:BE}
In the setting of \cref{prop:tangent}, $pH-2tm\ge0$.
\end{corollary}

The second is the line of \cref{lem:line} below.

\subsection{The supporting lines at the ends}

\begin{lemma}\label{lem:ends}
For every $s\in\R$,
\[
  \tilde h(s)\ge h(1)\Bigl(1+\frac p2(s-1)\Bigr)\qquad\text{and}\qquad
  \tilde h(s)\ge h(-1)\Bigl(1-\frac p2(s+1)\Bigr).
\]
\end{lemma}

At the conformal point the slope of $h$ is therefore at most $\tfrac p2h(1)$. For $p\le2$ the
first inequality says that $F\ge F(I)\,B_p/B_p(I)$ when $\Phi$ is the profile of $F$: the line
$s\mapsto1+\tfrac p2(s-1)$ is, up to the factor $B_p(I)=2/p$, the function $h$ of $B_p$. This is
the extremality of the Burkholder function; cf.~\cite[Theorem~4.2]{GuerraKristensen2022}
and~\cite{AIPS2015}.
